\IfFileExists{stacks-project.cls}{%
\documentclass{stacks-project}
}{%
\documentclass{amsart}
}

% For dealing with references we use the comment environment
\usepackage{verbatim}
\newenvironment{reference}{\comment}{\endcomment}
%\newenvironment{reference}{}{}
\newenvironment{slogan}{\comment}{\endcomment}
\newenvironment{history}{\comment}{\endcomment}

% For commutative diagrams we use Xy-pic
\usepackage[all]{xy}

% We use 2cell for 2-commutative diagrams.
\xyoption{2cell}
\UseAllTwocells

% We use multicol for the list of chapters between chapters
\usepackage{multicol}

% This is generally recommended for better output
\usepackage{lmodern}
\usepackage[T1]{fontenc}

% For cross-file-references
\usepackage{xr-hyper}

% Package for hypertext links:
\usepackage{hyperref}

% For any local file, say "hello.tex" you want to link to please
% use \externaldocument[hello-]{hello}
\externaldocument[introduction-]{introduction}
\externaldocument[conventions-]{conventions}
\externaldocument[sets-]{sets}
\externaldocument[categories-]{categories}
\externaldocument[topology-]{topology}
\externaldocument[sheaves-]{sheaves}
\externaldocument[sites-]{sites}
\externaldocument[stacks-]{stacks}
\externaldocument[fields-]{fields}
\externaldocument[algebra-]{algebra}
\externaldocument[brauer-]{brauer}
\externaldocument[homology-]{homology}
\externaldocument[derived-]{derived}
\externaldocument[simplicial-]{simplicial}
\externaldocument[more-algebra-]{more-algebra}
\externaldocument[smoothing-]{smoothing}
\externaldocument[modules-]{modules}
\externaldocument[sites-modules-]{sites-modules}
\externaldocument[injectives-]{injectives}
\externaldocument[cohomology-]{cohomology}
\externaldocument[sites-cohomology-]{sites-cohomology}
\externaldocument[dga-]{dga}
\externaldocument[dpa-]{dpa}
\externaldocument[sdga-]{sdga}
\externaldocument[hypercovering-]{hypercovering}
\externaldocument[schemes-]{schemes}
\externaldocument[constructions-]{constructions}
\externaldocument[properties-]{properties}
\externaldocument[morphisms-]{morphisms}
\externaldocument[coherent-]{coherent}
\externaldocument[divisors-]{divisors}
\externaldocument[limits-]{limits}
\externaldocument[varieties-]{varieties}
\externaldocument[topologies-]{topologies}
\externaldocument[descent-]{descent}
\externaldocument[perfect-]{perfect}
\externaldocument[more-morphisms-]{more-morphisms}
\externaldocument[flat-]{flat}
\externaldocument[groupoids-]{groupoids}
\externaldocument[more-groupoids-]{more-groupoids}
\externaldocument[etale-]{etale}
\externaldocument[chow-]{chow}
\externaldocument[intersection-]{intersection}
\externaldocument[pic-]{pic}
\externaldocument[weil-]{weil}
\externaldocument[adequate-]{adequate}
\externaldocument[dualizing-]{dualizing}
\externaldocument[duality-]{duality}
\externaldocument[discriminant-]{discriminant}
\externaldocument[derham-]{derham}
\externaldocument[local-cohomology-]{local-cohomology}
\externaldocument[algebraization-]{algebraization}
\externaldocument[curves-]{curves}
\externaldocument[resolve-]{resolve}
\externaldocument[models-]{models}
\externaldocument[functors-]{functors}
\externaldocument[equiv-]{equiv}
\externaldocument[pione-]{pione}
\externaldocument[etale-cohomology-]{etale-cohomology}
\externaldocument[proetale-]{proetale}
\externaldocument[relative-cycles-]{relative-cycles}
\externaldocument[more-etale-]{more-etale}
\externaldocument[trace-]{trace}
\externaldocument[crystalline-]{crystalline}
\externaldocument[spaces-]{spaces}
\externaldocument[spaces-properties-]{spaces-properties}
\externaldocument[spaces-morphisms-]{spaces-morphisms}
\externaldocument[decent-spaces-]{decent-spaces}
\externaldocument[spaces-cohomology-]{spaces-cohomology}
\externaldocument[spaces-limits-]{spaces-limits}
\externaldocument[spaces-divisors-]{spaces-divisors}
\externaldocument[spaces-over-fields-]{spaces-over-fields}
\externaldocument[spaces-topologies-]{spaces-topologies}
\externaldocument[spaces-descent-]{spaces-descent}
\externaldocument[spaces-perfect-]{spaces-perfect}
\externaldocument[spaces-more-morphisms-]{spaces-more-morphisms}
\externaldocument[spaces-flat-]{spaces-flat}
\externaldocument[spaces-groupoids-]{spaces-groupoids}
\externaldocument[spaces-more-groupoids-]{spaces-more-groupoids}
\externaldocument[bootstrap-]{bootstrap}
\externaldocument[spaces-pushouts-]{spaces-pushouts}
\externaldocument[spaces-chow-]{spaces-chow}
\externaldocument[groupoids-quotients-]{groupoids-quotients}
\externaldocument[spaces-more-cohomology-]{spaces-more-cohomology}
\externaldocument[spaces-simplicial-]{spaces-simplicial}
\externaldocument[spaces-duality-]{spaces-duality}
\externaldocument[formal-spaces-]{formal-spaces}
\externaldocument[restricted-]{restricted}
\externaldocument[spaces-resolve-]{spaces-resolve}
\externaldocument[formal-defos-]{formal-defos}
\externaldocument[defos-]{defos}
\externaldocument[cotangent-]{cotangent}
\externaldocument[examples-defos-]{examples-defos}
\externaldocument[algebraic-]{algebraic}
\externaldocument[examples-stacks-]{examples-stacks}
\externaldocument[stacks-sheaves-]{stacks-sheaves}
\externaldocument[criteria-]{criteria}
\externaldocument[artin-]{artin}
\externaldocument[quot-]{quot}
\externaldocument[stacks-properties-]{stacks-properties}
\externaldocument[stacks-morphisms-]{stacks-morphisms}
\externaldocument[stacks-limits-]{stacks-limits}
\externaldocument[stacks-cohomology-]{stacks-cohomology}
\externaldocument[stacks-perfect-]{stacks-perfect}
\externaldocument[stacks-introduction-]{stacks-introduction}
\externaldocument[stacks-more-morphisms-]{stacks-more-morphisms}
\externaldocument[stacks-geometry-]{stacks-geometry}
\externaldocument[moduli-]{moduli}
\externaldocument[moduli-curves-]{moduli-curves}
\externaldocument[examples-]{examples}
\externaldocument[exercises-]{exercises}
\externaldocument[guide-]{guide}
\externaldocument[desirables-]{desirables}
\externaldocument[coding-]{coding}
\externaldocument[obsolete-]{obsolete}
\externaldocument[fdl-]{fdl}
\externaldocument[index-]{index}

% Theorem environments.
%
\theoremstyle{plain}
\newtheorem{theorem}[subsection]{Theorem}
\newtheorem{proposition}[subsection]{Proposition}
\newtheorem{lemma}[subsection]{Lemma}

\theoremstyle{definition}
\newtheorem{definition}[subsection]{Definition}
\newtheorem{example}[subsection]{Example}
\newtheorem{exercise}[subsection]{Exercise}
\newtheorem{situation}[subsection]{Situation}

\theoremstyle{remark}
\newtheorem{remark}[subsection]{Remark}
\newtheorem{remarks}[subsection]{Remarks}

\numberwithin{equation}{subsection}

% Macros
%
\def\lim{\mathop{\mathrm{lim}}\nolimits}
\def\colim{\mathop{\mathrm{colim}}\nolimits}
\def\Spec{\mathop{\mathrm{Spec}}}
\def\Hom{\mathop{\mathrm{Hom}}\nolimits}
\def\Ext{\mathop{\mathrm{Ext}}\nolimits}
\def\SheafHom{\mathop{\mathcal{H}\!\mathit{om}}\nolimits}
\def\SheafExt{\mathop{\mathcal{E}\!\mathit{xt}}\nolimits}
\def\Sch{\mathit{Sch}}
\def\Mor{\mathop{\mathrm{Mor}}\nolimits}
\def\Ob{\mathop{\mathrm{Ob}}\nolimits}
\def\Sh{\mathop{\mathit{Sh}}\nolimits}
\def\NL{\mathop{N\!L}\nolimits}
\def\CH{\mathop{\mathrm{CH}}\nolimits}
\def\proetale{{pro\text{-}\acute{e}tale}}
\def\etale{{\acute{e}tale}}
\def\QCoh{\mathit{QCoh}}
\def\Ker{\mathop{\mathrm{Ker}}}
\def\Im{\mathop{\mathrm{Im}}}
\def\Coker{\mathop{\mathrm{Coker}}}
\def\Coim{\mathop{\mathrm{Coim}}}

% Boxtimes
%
\DeclareMathSymbol{\boxtimes}{\mathbin}{AMSa}{"02}

%
% Macros for moduli stacks/spaces
%
\def\QCohstack{\mathcal{QC}\!\mathit{oh}}
\def\Cohstack{\mathcal{C}\!\mathit{oh}}
\def\Spacesstack{\mathcal{S}\!\mathit{paces}}
\def\Quotfunctor{\mathrm{Quot}}
\def\Hilbfunctor{\mathrm{Hilb}}
\def\Curvesstack{\mathcal{C}\!\mathit{urves}}
\def\Polarizedstack{\mathcal{P}\!\mathit{olarized}}
\def\Complexesstack{\mathcal{C}\!\mathit{omplexes}}
% \Pic is the operator that assigns to X its picard group, usage \Pic(X)
% \Picardstack_{X/B} denotes the Picard stack of X over B
% \Picardfunctor_{X/B} denotes the Picard functor of X over B
\def\Pic{\mathop{\mathrm{Pic}}\nolimits}
\def\Picardstack{\mathcal{P}\!\mathit{ic}}
\def\Picardfunctor{\mathrm{Pic}}
\def\Deformationcategory{\mathcal{D}\!\mathit{ef}}
\usepackage{cleveref}
\crefname{originalreference}{original reference}{original references}
\Crefname{originalreference}{Original reference}{Original references}
\makeatletter
\newlabel{definition-cartesian-derived-modules}{{0D9N}{source}{}{sourcepointer-1}{}}
\newlabel{definition-cartesian-derived-modules@cref}{{[originalreference][0][]0D9N}{[1][0][]source}{}{}{}}
\makeatother
\crefname{originalreference}{original reference}{original references}
\Crefname{originalreference}{Original reference}{Original references}
\makeatletter
\newlabel{lemma-cartesian-module-derived-from-simplicial}{{0D9T}{source}{}{sourcepointer-2}{}}
\newlabel{lemma-cartesian-module-derived-from-simplicial@cref}{{[originalreference][0][]0D9T}{[1][0][]source}{}{}{}}
\makeatother
\crefname{originalreference}{original reference}{original references}
\Crefname{originalreference}{Original reference}{Original references}
\makeatletter
\newlabel{lemma-hom-cartesian-modules-derived}{{0D9S}{source}{}{sourcepointer-3}{}}
\newlabel{lemma-hom-cartesian-modules-derived@cref}{{[originalreference][0][]0D9S}{[1][0][]source}{}{}{}}
\makeatother
\crefname{originalreference}{original reference}{original references}
\Crefname{originalreference}{Original reference}{Original references}
\makeatletter
\newlabel{lemma-nullity-cartesian-modules-derived}{{0D9R}{source}{}{sourcepointer-4}{}}
\newlabel{lemma-nullity-cartesian-modules-derived@cref}{{[originalreference][0][]0D9R}{[1][0][]source}{}{}{}}
\makeatother
\crefname{originalreference}{original reference}{original references}
\Crefname{originalreference}{Original reference}{Original references}
\makeatletter
\newlabel{remark-augmentation-ringed}{{0D9W}{source}{}{sourcepointer-5}{}}
\newlabel{remark-augmentation-ringed@cref}{{[originalreference][0][]0D9W}{[1][0][]source}{}{}{}}
\makeatother
\crefname{originalsection}{original section}{original sections}
\Crefname{originalsection}{Original section}{Original sections}
\makeatletter
\newlabel{section-semi-representable}{{09WK}{source}{}{sourcepointer-6}{}}
\newlabel{section-semi-representable@cref}{{[originalsection][0][]09WK}{[1][0][]source}{}{}{}}
\makeatother
\crefname{originalreference}{original reference}{original references}
\Crefname{originalreference}{Original reference}{Original references}
\makeatletter
\newlabel{situation-locally-given}{{0D67}{source}{}{sourcepointer-7}{}}
\newlabel{situation-locally-given@cref}{{[originalreference][0][]0D67}{[1][0][]source}{}{}{}}
\makeatother
\title{Existence and uniqueness of the cartesian derived glueing along simplicial semi-representable objects}
\author{The Stacks Project Authors}
\date{stacks-project:spaces-simplicial.tex:local-snapshot-20261012; GNU FDL 1.2 or later}
\begin{document}
\maketitle
\noindent\textit{Editorial note.} All delivered mathematical text is the author's own, copied from the pinned Stacks Project chapter file sources/stacks-spaces-simplicial/pinned-original-source.tex; no statement or proof step was rewritten, reordered or removed.\par\medskip
\noindent\textit{Editorial note.} The counted proof is the author's complete argument: the 'formal' product terminology is the author's (source lines 5607-5608), referring to the product decomposition of the derived categories from the semi-representable section; the vanishing Hom(E\_n[t], E\_n) = 0 for t > 0 follows from the assumed vanishing of negative self-exts (source lines 5614-5616); this lemma prepares the BBD glueing theorems delivered as cards L46160 (cech variant) and L46162 (bbd-glueing).\par\medskip
\noindent\textit{Editorial note.} The label lemma-prepare-bbd-glueing names the counted statement, and the proof environment immediately following it in the source is delivered verbatim as its proof.\par\medskip
\noindent\textit{Editorial note.} The counted statement is carried by the author's own lemma environment in the chapter source.\par\medskip
\noindent\textit{Allowed standard prerequisites (editorial).}\begin{itemize}
\item The locally-given situation and the negative-ext assumption. Same chapter, Situation \textbackslash{}label\{situation-locally-given\}. Statement lines 5589-5592: negative self-exts of E\_U in D(O\_\{u(U)\}) are zero; L is a simplicial object of SR(B) with K = u(L).
\item The ringed total site of the augmentation. Same chapter, Remark \textbackslash{}label\{remark-augmentation-ringed\}. Statement lines 5593-5596: ((C/K)\_total, O) is the ringed total site of the remark, with component ringed topoi O\_n on C/K\_n = C/u(L\_n).
\item The product decomposition of semi-representable topoi. Same chapter, Section \textbackslash{}label\{section-semi-representable\}. Proof lines 5602-5607: Sh(C/K\_n) = prod Sh(C/u(U\_\{n,i\})) and similarly for modules and derived categories, compatibly with the morphisms of the simplicial object K.
\item Cartesian derived objects from simplicial systems, and uniqueness. Same chapter, \textbackslash{}label\{lemma-cartesian-module-derived-from-simplicial\}, \textbackslash{}label\{lemma-nullity-cartesian-modules-derived\} and \textbackslash{}label\{lemma-hom-cartesian-modules-derived\}. Proof lines 5615-5627: the simplicial system with Hom(E\_n[t], E\_n) = 0 for t > 0 produces the cartesian object E; uniqueness up to unique isomorphism follows from the nullity and Hom lemmas.
\end{itemize}\medskip
\ifdefined\hypertarget\hypertarget{sourcepointer-1}{}\fi
\noindent\textit{Source reference.} Original reference 0D9N: sources/stacks-spaces-simplicial/pinned-original-source.tex, source lines 3357--3357. Prerequisite pointer invoked at proof line(s) 5618; it names a standard fact of the same chapter that the counted argument uses, and every step of the counted proof built on it is retained verbatim in the delivered body.\par\medskip
\ifdefined\hypertarget\hypertarget{sourcepointer-2}{}\fi
\noindent\textit{Source reference.} Original reference 0D9T: sources/stacks-spaces-simplicial/pinned-original-source.tex, source lines 3539--3539. Prerequisite pointer invoked at proof line(s) 5622; it names a standard fact of the same chapter that the counted argument uses, and every step of the counted proof built on it is retained verbatim in the delivered body.\par\medskip
\ifdefined\hypertarget\hypertarget{sourcepointer-3}{}\fi
\noindent\textit{Source reference.} Original reference 0D9S: sources/stacks-spaces-simplicial/pinned-original-source.tex, source lines 3516--3516. Prerequisite pointer invoked at proof line(s) 5626; it names a standard fact of the same chapter that the counted argument uses, and every step of the counted proof built on it is retained verbatim in the delivered body.\par\medskip
\ifdefined\hypertarget\hypertarget{sourcepointer-4}{}\fi
\noindent\textit{Source reference.} Original reference 0D9R: sources/stacks-spaces-simplicial/pinned-original-source.tex, source lines 3493--3493. Prerequisite pointer invoked at proof line(s) 5625; it names a standard fact of the same chapter that the counted argument uses, and every step of the counted proof built on it is retained verbatim in the delivered body.\par\medskip
\ifdefined\hypertarget\hypertarget{sourcepointer-5}{}\fi
\noindent\textit{Source reference.} Original reference 0D9W: sources/stacks-spaces-simplicial/pinned-original-source.tex, source lines 4073--4073. Prerequisite pointer invoked at proof line(s) 5594; it names a standard fact of the same chapter that the counted argument uses, and every step of the counted proof built on it is retained verbatim in the delivered body.\par\medskip
\ifdefined\hypertarget\hypertarget{sourcepointer-6}{}\fi
\noindent\textit{Source reference.} Original section 09WK: sources/stacks-spaces-simplicial/pinned-original-source.tex, source lines 3578--3578. Prerequisite pointer invoked at proof line(s) 5609; it names a standard fact of the same chapter that the counted argument uses, and every step of the counted proof built on it is retained verbatim in the delivered body.\par\medskip
\ifdefined\hypertarget\hypertarget{sourcepointer-7}{}\fi
\noindent\textit{Source reference.} Original reference 0D67: sources/stacks-spaces-simplicial/pinned-original-source.tex, source lines 5552--5552. Prerequisite pointer invoked at proof line(s) 5589; it names a standard fact of the same chapter that the counted argument uses, and every step of the counted proof built on it is retained verbatim in the delivered body.\par\medskip
\begin{lemma}
\label{lemma-prepare-bbd-glueing}
In Situation \ref{situation-locally-given}.
Assume negative self-exts of $E_U$ in $D(\mathcal{O}_{u(U)})$ are zero.
Let $L$ be a simplicial object of $\text{SR}(\mathcal{B})$.
Consider the simplicial object $K = u(L)$ of $\text{SR}(\mathcal{C})$
and let $((\mathcal{C}/K)_{total}, \mathcal{O})$ be as in
Remark \ref{remark-augmentation-ringed}.
There exists a cartesian object $E$ of $D(\mathcal{O})$
such that writing $L_n = \{U_{n, i}\}_{i \in I_n}$
the restriction of $E$ to $D(\mathcal{O}_{\mathcal{C}/u(U_{n, i})})$
is $E_{U_{n, i}}$ compatibly (see proof for details).
Moreover, $E$ is unique up to unique isomorphism.
\end{lemma}

\begin{proof}
Recall that
$\Sh(\mathcal{C}/K_n) = \prod_{i \in I_n} \Sh(\mathcal{C}/u(U_{n, i}))$
and similarly for the categories of modules. This product decomposition
is also inherited by the derived categories of sheaves of modules.
Moreover, this product decomposition is compatible with
the morphisms in the simplicial semi-representable object $K$.
See Section \ref{section-semi-representable}.
Hence we can set $E_n = \prod_{i \in I_n} E_{U_{n, i}}$
(``formal'' product) in $D(\mathcal{O}_n)$.
Taking (formal) products of the maps $\rho_a$ of
Situation \ref{situation-locally-given}
we obtain isomorphisms $E_\varphi : f_\varphi^*E_n \to E_m$.
The assumption about compostions of the maps $\rho_a$
immediately implies that $(E_n, E_\varphi)$
defines a simplicial system of the derived category of modules
as in Definition \ref{definition-cartesian-derived-modules}.
The vanishing of negative exts assumed in the lemma implies that
$\Hom(E_n[t], E_n) = 0$ for $n \geq 0$ and $t > 0$.
Thus by
Lemma \ref{lemma-cartesian-module-derived-from-simplicial}
we obtain $E$.
Uniqueness up to unique isomorphism follows from
Lemmas \ref{lemma-nullity-cartesian-modules-derived} and
\ref{lemma-hom-cartesian-modules-derived}.
\end{proof}

\end{document}
